% =============================================================
%  要买的东西
% =============================================================
\Chapter{开工前：要买什么、翻什么、印什么}
\phantomsection\label{chap:supplies}

\begin{unitinfo}
二年级上册凡是“摆一摆、量一量、剪一剪、买一买”的地方，都得有东西上手。
没东西，小豆角就只能看爸爸拿笔在纸上画——那跟看视频没区别，学不进去。\\
\textbf{一次买齐，整学期不用再跑文具店。}清单不长，看着像玩具，其实是这学期最要紧的一套家伙。
\end{unitinfo}

\subsection{为什么非要买小棒}

一包小棒，看着像玩具，其实是这学期最核心的一件教具。

因为\textbf{十进制这件事，只能用“捆”来解释}。爸爸说“满十进一”，
小豆角的脑子里没有画面；爸爸把 13 根散的让小豆角自己捆成一捆，
小豆角手上就有了一个动作：\textbf{10 根变成 1 捆}。
这个动作重复十几次，小豆角才真的知道“十位”是什么东西。

同样的道理：
\begin{itemize}
  \item \textbf{计数器}——解释“100 减 48 借不到怎么办”，比小棒更直观（珠子一挪就看得见）。
  \item \textbf{直尺}——不摸尺子，“1 厘米”永远是个词。
  \item \textbf{硬币纸币}——不算钱，“1 元 = 10 角”就是背书。
\end{itemize}

\subsection{要买的（16 样）}

\begin{center}
\small
\begin{tabularx}{\linewidth}{@{}p{0.5cm}p{3.2cm}p{4.8cm}p{1.0cm}X@{}}
\toprule
\textbf{序} & \textbf{买什么} & \textbf{怎么挑} & \textbf{几份} & \textbf{用在哪} \\
\midrule
1 & 小棒 & 木质或塑料计数棒，一包 100 根。\textbf{重点：要能捆}，别买成一整块塑料板 &
      2 包 & 全书主力，一、三、四、五、七、八单元 \\
2 & 计数器 & 五档（个十百千万），每档 10 颗珠 & 1 个 & 第一单元进位退位、第二单元读数 \\
3 & 透明直尺 15\,cm & 刻度清楚，0 刻度在最边上 & 2 把 & 第二单元，两把能接起来量长的 \\
4 & 透明直尺 30\,cm & 同上 & 1 把 & 第二单元量课桌、书本 \\
5 & 软尺 & 1.5 米裁缝尺 & 1 条 & 第二单元量黑板、身高 \\
6 & 正方形手工纸 & 15\,cm$\times$15\,cm，100 张。\textbf{用量很大} & 1 包 & 第六单元剪对称 \\
7 & A4 彩色卡纸 & 20 张，多色 & 1 包 & 做陀螺、竹蜻蜓、价格牌 \\
8 & 儿童剪刀 & 圆头 & 1 把 & 第六单元 \\
9 & 固体胶 & 1 支 & 1 支 & 贴记录、做手工 \\
10 & 塑料吸管 & 普通饮料吸管，50 支 & 1 包 & 第六单元竹蜻蜓 \\
11 & 竹签／牙签 & 一头尖 & 1 盒 & 第六单元陀螺中轴 \\
12 & 图钉 & 1 盒 & 1 盒 & 第六单元风车、陀螺 \\
13 & 骰子 & 六面，1--6 点 & 2 颗 & 第五单元火星探险、口算游戏 \\
14 & 棋子 & 跳棋那种小圆片 & 1 副 & 第五单元棋盘 \\
15 & 贴纸 & 100 张装 & 1 包 & 全书星星奖励 \\
16 & 打印复印 & A4 约 30 页，其中 4 页用 160\,g 卡纸 & 1 次 & 附录全部 \\
\bottomrule
\end{tabularx}
\end{center}

\subsection{想再加点（可选）}

\begin{center}
\small
\begin{tabularx}{\linewidth}{@{}p{0.5cm}p{3.2cm}p{9.0cm}p{1.0cm}@{}}
\toprule
\textbf{序} & \textbf{买什么} & \textbf{为什么值} & \textbf{几份} \\
\midrule
17 & 磁性小白板 45$\times$60 & 讲竖式、画图能反复擦，比纸省事太多。
      不想买就\textbf{把 A4 白纸塞进透明文件袋}，用白板笔写，效果差不多 & 1 块 \\
18 & 收纳盒 & 3 个带盖的，分“计数／手工／纸品”，贴标签 & 3 个 \\
19 & 扑克牌 & 抽牌比大小、凑十、口诀对练，一学期都用得上 & 1 副 \\
\bottomrule
\end{tabularx}
\end{center}

\subsection{家里翻出来的}

\begin{center}
\small
\begin{tabularx}{\linewidth}{@{}p{3.4cm}Xp{3.0cm}@{}}
\toprule
\textbf{翻什么} & \textbf{拿来干嘛} & \textbf{哪个单元} \\
\midrule
1 米长的绳子／鞋带 & 剪成正好 1 米，去比家里哪些东西比 1 米长 & 第二单元 \\
豆子／玻璃珠／弹弓子弹 & 当计数物，比小棒好摆成堆 & 一、三、五单元 \\
纸杯 6 个 & 当盘子、袋子、笼子，做平均分 & 第五单元 \\
小饼干／糖果一包 & 分东西的真实版，\textbf{分完能吃} & 第五单元 \\
快递纸箱 & 拆开做华容道底板、陀螺纸板 & 第六单元 \\
旧杂志／报纸 & 练剪直线曲线，剪坏了不心疼 & 第六单元 \\
皮筋 10 根 & 把小棒捆成十根一捆，进位全靠它 & 第一单元 \\
手机 & 指南针（认方向）、计时、拍下小豆角的作品 & 全单元 \\
卷尺 & 量身高、量家具 & 第二单元 \\
透明文件袋 & 塞白纸当白板用 & 全单元 \\
\bottomrule
\end{tabularx}
\end{center}

\subsection{要印的（本书附录，约 30 页）}

\begin{center}
\small
\begin{tabularx}{\linewidth}{@{}p{1.5cm}p{5.0cm}p{1.2cm}X@{}}
\toprule
\textbf{在哪} & \textbf{印什么} & \textbf{几页} & \textbf{怎么印} \\
\midrule
附录 A & 口算卡（加减／乘法／除法，剪成小卡） & 8 页 & \textbf{160\,g 卡纸}，单面，剪开 \\
附录 B & 口诀表（完整＋填空）、百数表、数位表、方格纸 & 6 页 & 普通纸；口诀表和百数表多印两份 \\
附录 C & 模拟人民币 & 2 页 & \textbf{160\,g 卡纸}，剪开 \\
附录 D & 打卡表、错题表、星星表、速查卡 & 3 页 & 普通纸，速查卡贴冰箱 \\
第五单元内 & 火星探险棋盘 & 1 页 & 普通纸或卡纸 \\
第六单元内 & 华容道棋盘 & 1 页 & \textbf{卡纸}，剪成小块 \\
\bottomrule
\end{tabularx}
\end{center}

\begin{tipbox}
\textbf{打印一定要选“实际大小／100\%”}，千万别勾“适应页面”。
缩过的手工纸和人民币尺寸就不准了——第五单元买东西会算错钱。
\end{tipbox}

\subsection{每个单元额外要备的（前一天晚上花两分钟）}

\begin{center}
\small
\begin{tabularx}{\linewidth}{@{}p{2.2cm}X@{}}
\toprule
\textbf{单元} & \textbf{这一天额外要拿出来的} \\
\midrule

第一单元 & 小棒＋皮筋（捆十用）、计数器、骰子、方格竖式纸（附录 B） \\
第二单元 & 15\,cm 和 30\,cm 直尺、软尺、1 米绳子、
\textbf{一把“坏尺子”}（拿旧尺把 0 那头用胶带贴住，让 0 看不见） \\
第三单元 & 小棒、玻璃珠、弹弓子弹、点子图（附录 B） \\
第四单元 & 小棒、百数表、口诀卡（附录 A 的乘法卡先只用 1--5 部分） \\
第五单元 & 小棒 30 根、纸杯 6 个、一包小饼干、骰子和棋子、棋盘（本章末页） \\
第六单元 & 手工纸、卡纸、剪刀、吸管、竹签、图钉、快递纸箱 \\
第七单元 & 百数表、口诀填空表、\textbf{小豆角自己的十根手指}（9 的口诀要用） \\
第八单元 & 小棒、附录 A／B 的练习纸 \\
综合实践 & 手机指南针、A3 纸、模拟人民币、家里小物件 \\
\bottomrule
\end{tabularx}
\end{center}

\subsection{装盒，五分钟搞定}

\begin{center}
\small
\begin{tabularx}{\linewidth}{@{}p{2.4cm}p{1.0cm}X@{}}
\toprule
\textbf{盒子} & \textbf{贴的标签} & \textbf{装这些} \\
\midrule
盒 1 & 计数 & 小棒两包、皮筋、豆子、骰子、棋子、扑克牌 \\
盒 2 & 手工 & 剪刀、固体胶、手工纸、卡纸、吸管、竹签、图钉、1 米绳子 \\
盒 3 & 纸品 & 口算卡、口诀表、百数表、数位表、人民币、记录表、方格纸 \\
\bottomrule
\end{tabularx}
\end{center}

\noindent 盒 3 里再拿文件袋分一下（卡一袋、表一袋、钱一袋）。
上课前三十秒能拿全，才不会因为“找东西”把二十分钟耗掉。

\subsection{直接抄走，下单去}

\begin{tipbox}
小棒 ×2 包\quad 计数器（五档）×1\quad 透明直尺 15\,cm ×2\quad 透明直尺 30\,cm ×1\\
软尺 1.5\,m ×1\quad 正方形手工纸 15\,cm ×1 包\quad A4 彩色卡纸 ×1 包\quad 儿童剪刀 ×1\\
固体胶 ×1\quad 塑料吸管 ×1 包\quad 竹签 ×1 盒\quad 图钉 ×1 盒\\
骰子 ×2\quad 跳棋棋子 ×1 副\quad 奖励贴纸 ×1 包\\
（可选）磁性白板 ×1\quad 收纳盒 ×3\quad 扑克牌 ×1
\end{tipbox}

\noindent \textbf{优先级}：小棒、计数器、直尺、手工纸、剪刀——这五样\textbf{第一周就得到位}，
因为第一、二单元立刻要用。剩下的边用边补，不着急。
